\clearpage
\subsection{RQF-05 — Bitácora apícola con calendario}

\renewcommand{\arraystretch}{1.3}

\noindent
\begin{tabular}{|p{5cm}|p{10.5cm}|}
\hline
\rowcolor{secVerde}\multicolumn{2}{|l|}{\color{white}\textbf{1. INFORMACIÓN DEL CLIENTE}} \\
\hline
Nombre de la empresa  & UTEQ \\ \hline
Nombre del Contacto   & Angélica Garduño Bustamante \\ \hline
Fecha de Elaboración  & 01/06/2026 \\ \hline
\end{tabular}

\vspace{0.3cm}

\noindent
\begin{tabular}{|p{5cm}|p{10.5cm}|}
\hline
\rowcolor{secVerde}\multicolumn{2}{|l|}{\color{white}\textbf{2. INFORMACIÓN GENERAL}} \\
\hline
Nombre del Requerimiento     & RQF 05. Bitácora apícola con calendario \\ \hline
Tipo de Requerimiento        & Funcional \\ \hline
Fecha requerida de solución  & 30/07/2026 \\ \hline
Prioridad                    & Alta \\ \hline
\end{tabular}

\vspace{0.3cm}

\noindent
\begin{tabular}{|p{4.5cm}|p{11cm}|}
\hline
\rowcolor{secVerde}\multicolumn{2}{|l|}{\color{white}\textbf{3. DESCRIPCIÓN DEL REQUERIMIENTO}} \\
\hline
Objetivo del Requerimiento &
Implementar una bitácora apícola digital con calendario interactivo que permita al apicultor registrar, consultar, editar y eliminar eventos asociados a sus colmenas y apiarios, sustituyendo la libreta de campo tradicional. \\
\hline
Beneficio esperado &
Centralización del historial de inspecciones, tratamientos, cosechas y cambios de reina en un solo lugar accesible desde el dispositivo móvil. Mejora la trazabilidad de las labores apícolas y facilita la toma de decisiones informadas. \\
\hline
Alcance &
- Crear nueva pantalla \textit{Bitácora} accesible desde el \textit{drawer} del usuario. \par
- Mostrar un calendario mensual con dots de color en los días con eventos. \par
- CRUD completo de eventos clasificados por tipo. \par
- Acceso rápido desde la pantalla de detalle de cada colmena. \par
- Crear tabla \texttt{bitacora} en PostgreSQL y endpoints REST asociados. \\
\hline
Ubicación & \url{https://github.com/MarkG777/AbejaNet} \\ \hline
Descripción de Usuarios y Responsables &
\textbf{Usuarios:} \par
\textbf{Apicultor (Cliente):} registra y consulta eventos sobre el manejo de sus colmenas. \par
\vspace{0.1cm}
\textbf{Responsables:} \par
\textbf{App Móvil:} ofrece la interfaz de calendario y los formularios de evento. \par
\textbf{API Backend:} expone los endpoints \texttt{/api/bitacora} para CRUD. \par
\textbf{Base de Datos PostgreSQL:} almacena los eventos en la tabla \texttt{bitacora} con relaciones a \texttt{usuarios}, \texttt{colmenas} y \texttt{apiarios}. \\
\hline
Descripción funcional &
1. El apicultor accede a \textit{Bitácora} desde el \textit{drawer} del usuario. \par
2. Visualiza un calendario mensual con dots de color en los días con eventos registrados. \par
3. Selecciona un día y consulta el listado de eventos de esa fecha. \par
4. Mediante un botón flotante (FAB) crea un nuevo evento eligiendo tipo (Inspección, Cosecha, Tratamiento, Cambio de reina, División/Unión, Alimentación o Nota libre), colmena/apiario asociado, título y comentario. \par
5. Puede editar o eliminar eventos existentes. \par
6. Desde la pantalla de detalle de una colmena, dispone de un acceso rápido \textit{``Registrar evento''} que pre-llena la colmena en el formulario. \par
7. Puede filtrar el calendario por colmena o apiario específico. \\
\hline
Restricciones &
- Cada evento debe estar asociado al menos a una colmena o a un apiario. \par
- Los comentarios libres tienen un límite de 1000 caracteres. \par
- Solo el usuario propietario puede ver y modificar sus propios eventos. \\
\hline
Precedencias & RQF 02. Registro de clientes \\ \hline
Documentación requerida & Manual de usuario con guía de uso de la bitácora y descripción de los tipos de evento. \\ \hline
Observaciones & Se introduce la librería \texttt{react-native-calendars} y se requiere actualización del esquema de base de datos. \\ \hline
\end{tabular}

\vspace{0.3cm}

\noindent
\begin{tabular}{|p{5.1cm}|p{5.1cm}|p{5.1cm}|}
\hline
\rowcolor{secVerde}\multicolumn{3}{|l|}{\color{white}\textbf{4. CONDICIONES DE PRUEBA}} \\
\hline
\rowcolor{cellGris}\textbf{Condición} & \textbf{Validación} & \textbf{Resultado} \\
\hline
El apicultor crea un nuevo evento de tipo \textit{Inspección} desde la pantalla de bitácora. &
El sistema debe persistir el evento en el backend y reflejarlo en el calendario. &
El día correspondiente muestra un dot de color y al seleccionarlo aparece el evento creado. \\
\hline
El apicultor accede a \textit{``Registrar evento''} desde el detalle de una colmena. &
El formulario debe abrirse con la colmena pre-seleccionada en el campo correspondiente. &
El evento se guarda asociado correctamente a esa colmena sin requerir selección manual. \\
\hline
\end{tabular}
